\section{What an integral measures}

Integration begins with a question that is complementary to differentiation. Differentiation asks how a function changes locally. Integration asks how quantities accumulate.

There are two closely related ways to meet integrals. An indefinite integral asks for a function whose derivative is known. A definite integral asks for accumulated signed area or total change over an interval. Both ideas are important, and both will appear in this chapter.

\subsection*{Undoing differentiation}

Suppose we know that
\[
F'(x)=f(x).
\]
Then \(F\) is called an antiderivative of \(f\). Finding such a function is one meaning of integration.

For example, since
\[
\frac{d}{dx}x^2=2x,
\]
we say that \(x^2\) is an antiderivative of \(2x\).

But it is not the only one. For every constant \(C\),
\[
\frac{d}{dx}(x^2+C)=2x.
\]
The derivative of a constant is zero, so differentiation loses constant information.

\begin{remarkbox}
Integration as antidifferentiation recovers a function from its derivative, but only up to an additive constant.
\end{remarkbox}

\subsection*{Accumulation}

The second meaning of integration is accumulation. If a quantity changes continuously, an integral can measure the total accumulated amount over an interval.

For a non-negative function, the definite integral
\[
\int_a^b f(x)\,dx
\]
can be read geometrically as the area under the graph of \(f\) from \(x=a\) to \(x=b\). If the function takes negative values, the integral measures signed area: regions below the horizontal axis contribute negatively.

This interpretation will become more precise later. For now, the important point is that integration adds up infinitely many small contributions.

\subsection*{The notation}

The symbol
\[
\int f(x)\,dx
\]
denotes an indefinite integral. It asks for all antiderivatives of \(f\).

The symbol
\[
\int_a^b f(x)\,dx
\]
denotes a definite integral. It has limits of integration \(a\) and \(b\), and its result is a number.

The expression \(dx\) reminds us of the variable of integration. It also hints at the idea of summing small pieces along the \(x\)-axis.

\begin{examplebox}
The expression
\[
\int 2x\,dx
\]
asks for functions whose derivative is \(2x\). The answer is
\[
x^2+C.
\]
The expression
\[
\int_0^3 2x\,dx
\]
asks for a number: the accumulated value of \(2x\) from \(0\) to \(3\). Later we will compute it as
\[
\left[x^2\right]_0^3=9.
\]
\end{examplebox}

\begin{summarybox}
When reading an integral, ask:
\begin{itemize}
\item Is the integral indefinite or definite?
\item Am I looking for a family of antiderivatives or a number?
\item What is the variable of integration?
\item Is the integral being used to undo differentiation or to measure accumulation?
\item What role does the constant of integration play?
\end{itemize}
\end{summarybox}

The integral is not only a symbol for a calculation. It is a language for recovering functions and measuring accumulated change.